\documentclass[11pt,a4paper]{article}
\usepackage[italian]{babel}
\usepackage[margin=2.5cm]{geometry}
\usepackage{parskip}
\usepackage{amsmath, amssymb, amsthm}
\usepackage{booktabs}
\usepackage{array}
\usepackage{xcolor}
\usepackage{needspace}
\usepackage{enumitem}
\usepackage{listings}
\usepackage{tikz}
\usepackage[most]{tcolorbox}
\usepackage[hidelinks]{hyperref}
\usetikzlibrary{decorations.pathreplacing, shapes.geometric, arrows.meta, positioning}

% Niente vedove/orfane: i paragrafi non lasciano righe isolate a fine/inizio pagina
\widowpenalty=10000
\clubpenalty=10000

% Elenchi più compatti
\setlist{itemsep=2pt, parsep=0pt, topsep=4pt, partopsep=0pt}

% --- Riquadri con bordo nero, senza numero (si spezzano tra due pagine) ---
% Uso: \begin{definizione} ... \end{definizione}
%      \begin{teorema}[Nome] ... \end{teorema}  (il nome è facoltativo)
\tcbset{nota/.style={enhanced jigsaw, breakable, colback=white,
  colframe=black, boxrule=0.5pt, sharp corners}}
\NewTColorBox{definizione}{}{nota,
  before upper={\textbf{Definizione.}\ }}
\NewTColorBox{teorema}{o}{nota,
  before upper={\textbf{Teorema\IfValueT{#1}{ (#1)}.}\ }}
\NewTColorBox{proposizione}{o}{nota,
  before upper={\textbf{Proposizione\IfValueT{#1}{ (#1)}.}\ }}
\NewTColorBox{proprieta}{o}{nota,
  before upper={\textbf{Proprietà\IfValueT{#1}{ (#1)}.}\ }}
\NewTColorBox{assioma}{o}{nota, boxrule=1pt,
  before upper={\textbf{Assioma\IfValueT{#1}{ (#1)}.}\ }}

% --- Ambienti semplici (senza riquadro) ---
% Il titolo sta in cima al blocco, anche se il contenuto inizia con un riquadro o
% con codice; e non resta mai da solo in fondo a una pagina.
% Uso: \begin{esempio}[Titolo facoltativo] ... \end{esempio}
\newcommand{\newrunin}[2]{%
  \NewDocumentEnvironment{#1}{o}{%
    \par\Needspace{4\baselineskip}\addvspace{\medskipamount}\noindent
    \textbf{#2}\IfValueT{##1}{ (##1)}\textbf{.}\ \ignorespaces}%
  {\par\addvspace{\medskipamount}}}
\newrunin{esempio}{Esempio}
\newrunin{esercizio}{Esercizio}
\newrunin{osservazione}{Osservazione}

% V e F nelle tavole di verità
\newcommand{\V}{\textsf{V}}
\newcommand{\F}{\textsf{F}}

% --- Codice C ---
% Uso: \begin{codice} ... \end{codice}      codice C numerato
%      \begin{comando} ... \end{comando}    comandi da terminale
%      \begin{risultato} ... \end{risultato} output di un programma
%      \lstinline|printf("ciao");|           codice dentro al testo
\lstdefinestyle{c}{
  language=C,
  basicstyle=\ttfamily\small,
  keywordstyle=\color{blue}\bfseries,
  commentstyle=\color{gray}\itshape,
  stringstyle=\color{red!70!black},
  numbers=left,
  numberstyle=\tiny\color{gray},
  numbersep=8pt,
  columns=flexible,
  keepspaces=true,
  breaklines=true,
  showstringspaces=false,
  tabsize=4,
  morekeywords={include, define},
  % lettere accentate nei commenti
  literate={à}{{\`a}}1 {è}{{\`e}}1 {é}{{\'e}}1 {ì}{{\`i}}1
           {ò}{{\`o}}1 {ù}{{\`u}}1 {È}{{\`E}}1 {À}{{\`A}}1
}
\lstset{style=c}
\newtcblisting{codice}{listing only, nota, left=6mm,
  listing options={style=c}}
\newtcblisting{comando}{listing only, nota,
  listing options={basicstyle=\ttfamily\small, numbers=none, language={},
    columns=flexible, keepspaces=true, breaklines=true}}
\newtcblisting{risultato}{listing only, enhanced jigsaw, breakable,
  colback=black!5, colframe=black, boxrule=0.5pt, sharp corners,
  listing options={basicstyle=\ttfamily\small, numbers=none, language={},
    columns=flexible, keepspaces=true, breaklines=true}}

% --- Diagrammi di flusso (simboli standard) ---
\tikzset{
  terminale/.style = {ellipse, draw, minimum width=2.5cm, minimum height=0.9cm},
  io/.style        = {trapezium, trapezium left angle=70, trapezium right angle=110,
                      draw, minimum width=2.5cm, minimum height=0.9cm},
  processo/.style  = {rectangle, draw, minimum width=2.5cm, minimum height=0.9cm},
  decisione/.style = {diamond, draw, aspect=2, minimum width=2.5cm},
  freccia/.style   = {thick, -{Stealth}}
}

\title{Interi, divisibilità e divisione con resto}
\author{Marco Cerbella}
\date{Lezione 11 -- 26 ottobre 2026}

\begin{document}
\maketitle

\section{Gli interi come struttura algebrica}

$\mathbb{Z} = \{0, \pm 1, \pm 2, \pm 3, \dots\}$ con la somma e il prodotto di interi.
\begin{itemize}
  \item $(\mathbb{Z}, +)$ è un gruppo abeliano: $0$ elemento neutro, l'opposto di $k \in \mathbb{Z}$ è $-k \in \mathbb{Z}$.
  \item $(\mathbb{Z}, \cdot)$ è un semigruppo commutativo, con $+1$ elemento neutro per $\cdot$ ($1$ è l'unità); $+1$ e $-1$ sono invertibili per $\cdot$.
  \item Valgono le distributive.
\end{itemize}
Quindi $(\mathbb{Z}, +, \cdot)$ è un anello commutativo con unità. Vale la legge di annullamento del prodotto: è un dominio di integrità. In particolare, se $k \neq 0$ e $k\cdot x = k\cdot y \Rightarrow x = y$ ($k, x, y \in \mathbb{Z}$): vale la legge di cancellazione del prodotto. Indicheremo $(\mathbb{Z}, +, \cdot)$ semplicemente con $\mathbb{Z}$.

\section{Divisibilità}

\begin{definizione}
Siano $a, b \in \mathbb{Z}$. Si dice che $a$ \emph{divide} $b$, e si scrive $a \mid b$, se $\exists k \in \mathbb{Z}$ tale che $b = k\cdot a$. Se $a$ non divide $b$ ($\nexists k$ tale che $b = ka$) si scrive $a \nmid b$.
\end{definizione}

\begin{esempio}
$5 \nmid 7$; \quad $3 \mid 33$ perché $3\cdot 11 = 33$; \quad $-4 \mid 20$ perché $\exists k = -5$ tale che $(-4)\cdot(-5) = 20$; \quad $7 \mid -21$ perché $\exists -3 \in \mathbb{Z}$ tale che $7\cdot(-3) = -21$.
\end{esempio}

\begin{osservazione}
\begin{itemize}
  \item $\forall a \in \mathbb{Z}$: $a \mid 0$, perché $0\cdot a = 0$.
  \item $\forall a \in \mathbb{Z}$: $1 \mid a$, perché $a\cdot 1 = a$.
  \item $0 \mid b \Rightarrow \exists k$ tale che $0\cdot k = b \Rightarrow b = 0$: cioè $0$ divide solo zero.
  \item $\forall a$: $a \mid a$, perché $a\cdot 1 = a$.
  \item La relazione $\mid$ su $\mathbb{N}$ è una relazione d'ordine parziale. Su $\mathbb{Z}$ non è una relazione d'ordine perché non è antisimmetrica: se $a \mid b \wedge b \mid a$ allora $a = b$ oppure $a = -b$ (ad esempio $5 \mid -5 \wedge -5 \mid 5$ ma $-5 \neq 5$).
  \item $a \mid b \wedge a \mid c \Rightarrow a \mid (b + c)$; inoltre $a \mid (b + c) \wedge a \mid b \Rightarrow a \mid c$.
\end{itemize}
\end{osservazione}

\begin{proposizione}[Lemma]
Se $a \mid b$ e $a \mid c$ allora $a \mid (\alpha b + \beta c)$ per ogni $\alpha, \beta \in \mathbb{Z}$.
\end{proposizione}

\begin{proof}
$a \mid b$ significa $\exists k_1 \in \mathbb{Z}$ tale che $b = k_1\cdot a$; $a \mid c$ significa $\exists k_2 \in \mathbb{Z}$ tale che $c = k_2\cdot a$. Allora
\[
  \alpha b + \beta c = \alpha k_1 a + \beta k_2 a = (\alpha k_1 + \beta k_2)\cdot a = h\cdot a ,
\]
con $h = \alpha k_1 + \beta k_2 \in \mathbb{Z}$. Cioè $a \mid (\alpha b + \beta c)$.
\end{proof}

\section{Divisione con resto}

\begin{proposizione}
Siano $a, b \in \mathbb{Z}$, $b \neq 0$. Allora esistono e sono unici $q, r \in \mathbb{Z}$ tali che
\[
  a = q\cdot b + r, \qquad 0 \le r < |b| ,
\]
dove $q$ è il \emph{quoziente} e $r$ il \emph{resto}.
\end{proposizione}

\begin{proof}[Dimostrazione dell'esistenza]
Sia $T = \{k \in \mathbb{Z} \mid k = a - mb,\ m \in \mathbb{Z},\ k \ge 0\} \subseteq \mathbb{N}$.

$T \neq \emptyset$? Se $a \ge 0$, allora $a \in T$ (con $m = 0$). Se $a < 0$, esiste $\bar m \in \mathbb{Z}$ tale che $\bar m b \le a$ (ad esempio $\bar m = -|a|\,b$, perché $\bar m b = -|a|\,b^2 \le -|a| = a$), cioè $0 \le a - \bar m b$, cioè $a - \bar m b \in T$. Quindi $T \neq \emptyset$.

$\mathbb{N}$ è bene ordinato, quindi $T$ ha minimo. Sia $r = \min T$. Poiché $r \in T$, $r = a - qb$ per qualche $q \in \mathbb{Z}$, cioè $a = qb + r$, e $r \ge 0$. Resta da far vedere che $r < |b|$.

Per assurdo sia $r \ge |b|$, cioè $r = |b| + r_0$ con $r_0 \ge 0$. Allora
\[
  0 \le r_0 = r - |b| = a - qb - |b| = \begin{cases} a - b(q + 1) & b > 0 \\ a - b(q - 1) & b < 0 \end{cases}
\]
cioè $r_0 \in T$ e $r_0 = r - |b| < r$, contro la minimalità di $r$: assurdo. Dunque $r < |b|$.
\end{proof}

\begin{proof}[Dimostrazione dell'unicità]
Supponiamo che ci siano due coppie $q, r$ e $q', r'$ tali che $a = qb + r = q'b + r'$ con $0 \le r, r' < |b|$. Supponiamo $r \le r'$. Allora
\[
  0 \le r' - r = (a - q'b) - (a - qb) = b(q - q') .
\]
Passando ai valori assoluti:
\[
  0 \le |q - q'|\,|b| = |r' - r| = r' - r \le r' < |b| ,
\]
cioè $|q - q'|\,|b| < |b|$, cioè $|q - q'| < 1$, cioè $|q - q'| = 0$, cioè $q' = q$. Da $qb + r = q'b + r'$ segue allora $r' = r$.
\end{proof}

\begin{esempio}
\begin{itemize}
  \item $a = 15$, $b = 6$: $15 = 6\cdot 2 + 3$ ($q = 2$, $r = 3$).
  \item $a = 15$, $b = -6$: $15 = (-6)\cdot(-2) + 3$.
  \item $a = -15$, $b = 6$: $-15 = 6\cdot(-2) - 3 = 6\cdot(-2) - 3 + 6 - 6 = 6\cdot(-3) + 3$, cioè quoziente $-3$ e resto $3$.
  \item $a = -15$, $b = -6$: $-15 = (-6)\cdot 2 - 3 = (-6)\cdot 3 + 3$.
  \item $a = -24$, $b = 7$: $-24 = 7\cdot(-4) + 4$.
  \item $a = 24$, $b = -7$: $24 = (-7)\cdot(-3) + 3$.
\end{itemize}
\end{esempio}

\section{Classi di resto modulo 7}

$a \equiv b \bmod 7$ se $7 \mid a - b$, cioè se $\exists k \in \mathbb{Z}$ tale che $a - b = 7k$, cioè $a = 7k + b$. Le classi sono
\begin{align*}
  [0]_7 &= \{0, \pm 7, \pm 14, \dots\},\\
  [1]_7 &= \{1,\ 7 + 1,\ 14 + 1, \dots, 7k + 1, \dots, -7 + 1,\ -14 + 1, \dots\},\\
  [2]_7 &= \{2,\ 7 + 2, \dots, 7k + 2, \dots, -7 + 2,\ -14 + 2, \dots\},\\
  [-1]_7 &= \{-1 = (-1)\cdot 7 + 6,\ 6,\ 6 + 7k, \dots\} = [6]_7 .
\end{align*}
Le classi di congruenza modulo $n$ di $\mathbb{Z}_n = \mathbb{Z}/{\equiv_n}$ sono esattamente $n$, cioè i resti possibili della divisione per $n$: $[0], [1], \dots, [n-1]$.

\end{document}
